% !TEX program = lualatex
% Transparencias de Fundamentos de las Matemáticas I
% Clase 07: Conjuntos II

\documentclass[aspectratio=169,11pt]{beamer}

\usepackage{fontspec}
\usepackage[spanish,es-nodecimaldot]{babel}
\usepackage{amsmath,amssymb,mathtools}
\usepackage{booktabs}
\usepackage{csquotes}
\usepackage{xcolor}

\usetheme{Madrid}
\usecolortheme{default}
\setbeamertemplate{navigation symbols}{}
\setbeamertemplate{footline}[frame number]

\definecolor{FdMBlue}{RGB}{20,55,100}
\definecolor{FdMRed}{RGB}{130,30,30}
\definecolor{FdMGreen}{RGB}{25,100,60}

\setbeamercolor{structure}{fg=FdMBlue}
\setbeamercolor{title}{fg=white,bg=FdMBlue}
\setbeamercolor{frametitle}{fg=white,bg=FdMBlue}
\setbeamercolor{block title}{fg=white,bg=FdMBlue}
\setbeamercolor{block body}{bg=blue!3}
\setbeamercolor{alerted text}{fg=FdMRed}

\setmainfont{Latin Modern Roman}
\setsansfont{Latin Modern Sans}
\setmonofont{Latin Modern Mono}

\newcommand{\N}{\mathbb{N}}
\newcommand{\Z}{\mathbb{Z}}
\newcommand{\Q}{\mathbb{Q}}
\newcommand{\R}{\mathbb{R}}
\newcommand{\C}{\mathbb{C}}
\newcommand{\Pow}{\mathcal{P}}
\newcommand{\id}{\operatorname{id}}
\newcommand{\mcd}{\operatorname{mcd}}
\newcommand{\mcm}{\operatorname{mcm}}

\title[FdM I -- Clase 07]{Conjuntos II}
\subtitle{Producto cartesiano, partes, familias e identidades}
\author{Antonio Falcó Montesinos}
\institute{Universidad CEU Cardenal Herrera}
\date{Fundamentos de las Matemáticas I}

\begin{document}

\begin{frame}
  \titlepage
\end{frame}

\begin{frame}{Objetivos de la clase}
\begin{itemize}
  \item Usar producto cartesiano y pares ordenados.
  \item Comprender el conjunto de las partes.
  \item Demostrar identidades conjuntistas mediante doble inclusión.
\end{itemize}
\end{frame}

\begin{frame}{Mapa de la clase}
\tableofcontents
\end{frame}

\section{Construcciones}

\begin{frame}{Construcciones}
\begin{block}{Producto cartesiano}
\begin{itemize}
  \item \(A\times B\) es el conjunto de pares ordenados \((a,b)\) con \(a\in A\) y \(b\in B\).
  \item El orden importa: en general \(A\times B\neq B\times A\).
  \item Es la base para relaciones y aplicaciones.
\end{itemize}
\end{block}
\end{frame}

\begin{frame}{Construcciones}
\begin{block}{Conjunto de las partes}
\begin{itemize}
  \item \(\mathcal{P}(A)\) es el conjunto de todos los subconjuntos de \(A\).
  \item Si \(A=\{1,2\}\), entonces \(\mathcal{P}(A)=\{\varnothing,\{1\},\{2\},\{1,2\}\}\).
  \item Sus elementos son conjuntos.
\end{itemize}
\end{block}
\end{frame}

\begin{frame}{Construcciones}
\begin{block}{Familias de conjuntos}
\begin{itemize}
  \item Una familia de conjuntos es una colección indexada \((A_i)_{i\in I}\).
  \item Permite hablar de uniones e intersecciones de muchos conjuntos.
  \item El conjunto de índices debe estar especificado.
\end{itemize}
\end{block}
\end{frame}

\section{Demostrar identidades}

\begin{frame}{Demostrar identidades}
\begin{block}{Doble inclusión}
\begin{itemize}
  \item Para probar \(A=B\), se prueba \(A\subseteq B\) y \(B\subseteq A\).
  \item Cada inclusión se demuestra tomando un elemento arbitrario.
  \item La prueba debe traducir pertenencia en condiciones lógicas.
\end{itemize}
\end{block}
\end{frame}

\begin{frame}{Demostrar identidades}
\begin{block}{Ejemplo: De Morgan}
\begin{itemize}
  \item Para demostrar \((A\cup B)^c=A^c\cap B^c\), se toma \(x\) del lado izquierdo.
  \item La condición \(x\notin A\cup B\) equivale a \(x\notin A\) y \(x\notin B\).
  \item Eso equivale a \(x\in A^c\cap B^c\).
\end{itemize}
\end{block}
\end{frame}

\section{Profundización}

\begin{frame}{Profundización}
\begin{block}{Producto cartesiano en contexto}
\begin{itemize}
  \item Si \(A\) es un conjunto de estudiantes y \(B\) un conjunto de grupos, \(A\times B\) contiene pares estudiante-grupo.
  \item Un subconjunto de \(A\times B\) puede representar una relación de asignación.
  \item Las aplicaciones se definirán como un caso especial de esta idea.
\end{itemize}
\end{block}
\end{frame}

\begin{frame}{Profundización}
\begin{block}{Partes de un conjunto}
\begin{itemize}
  \item Si \(A\) tiene pocos elementos, \(\mathcal{P}(A)\) puede listarse explícitamente.
  \item El conjunto vacío y el conjunto total siempre pertenecen a \(\mathcal{P}(A)\).
  \item Cada elemento de \(\mathcal{P}(A)\) es a su vez un conjunto.
\end{itemize}
\end{block}
\end{frame}

\begin{frame}{Profundización}
\begin{block}{Doble inclusión: esquema}
\begin{itemize}
  \item Primera dirección: tomar \(x\) en el lado izquierdo y demostrar que está en el derecho.
  \item Segunda dirección: tomar \(x\) en el lado derecho y demostrar que está en el izquierdo.
  \item Al cerrar ambas direcciones se obtiene la igualdad de conjuntos.
\end{itemize}
\end{block}
\end{frame}

\begin{frame}{Cierre}
\begin{itemize}
  \item El producto cartesiano prepara relaciones y funciones.
  \item \(\mathcal{P}(A)\) contiene subconjuntos, no elementos sueltos.
  \item La doble inclusión es el método estándar para igualdades de conjuntos.
\end{itemize}
\end{frame}

\begin{frame}{Trabajo recomendado}
\begin{itemize}
  \item Releer las secciones correspondientes del manual.
  \item Identificar definiciones, símbolos nuevos y enunciados principales.
  \item Preparar dudas y ejemplos para el seminario asociado.
\end{itemize}
\end{frame}

\end{document}
